\section{Introduction}\label{sec:introduction84}
A measurement used as an unknown device returns a classical label, but
its discrimination need not be classical.  Correlations may be created
inside an acquisition, retained in an external receiver, or transmitted
coherently into a later acquisition.  These are distinct resources.
This article gives a local geometric classification with sharp allocation
laws, a general variational principle for reset decisions, and exact
finite experiments separating receiver storage from coherent acquisition.  These results concern the same classical-output
interface and retain the order of its labels.

\subsection{Local geometry and exact experiments}
An ordered measurement is a tuple $E=(E_1,\ldots,E_k)$ of positive
$d\times d$ matrices with $\sum_jE_j=I$.  Its channel is
$\mathcal M_E(\rho)=\sum_j\tr(E_j\rho)|j\rangle\langle j|$.
No post-measurement device system is available.  Known preparations,
ancillary systems and controls are allowed.  Trace norms and diamond
norms below are unhalved; success probabilities are ordinary probabilities.

The first result resolves the local geometry for a prescribed allocation.
Fix $E$, a nonzero Hermitian zero-sum tuple $H$ in the one-sided positive
tangent cone, and a remainder allowance $\Lambda\ge0$.  For
$\|R\|_\Sigma=\sum_j\|R_j\|_{\rm op}$, consider every legal $F$ with
\begin{equation}\label{eq:fronttube87}
 \|F-E-sH\|_\Sigma\le\Lambda s^2.
\end{equation}
Write $P_j=\operatorname{supp}E_j$, $Q_j=I-P_j$,
$J_j=Q_jH_jQ_j$, $\Gamma_E(H)=\frac i2\sum_j[P_j,H_j]$, and
$\mathcal W_E=\sum_jP_j\mathcal H_dP_j$, where $\mathcal H_d$ is the
real Hermitian matrix space.  The cone condition is $J_j\succeq0$.
For positive integer acquisition sizes $\boldsymbol n=(n_1,\ldots,n_\ell)$,
put $N=\sum_rn_r$ and $V=\sum_rn_r^2$.
Theorem~\ref{thm:allocation89} proves, respectively, the orders
\begin{equation}\label{eq:frontprofile89}
 \min\{1,\sqrt N s\},\qquad \min\{1,\sqrt V s\},\qquad \min\{1,Ns\}
\end{equation}
when $H\in\operatorname{Ran}\mathcal C_E$, when all $J_j=0$ and
$\Gamma_E(H)\notin\mathcal W_E$, and when some $J_j\ne0$.
The support-kernel identity shows these cases exhaust the nonzero cone.
The same orders hold for independent complete acquisition groups and
for prescribed-profile reset feedback.  At a reset cut the entire fresh
probe, reference and within-block memory is conditionally independent
of the receiver's entire old memory.  The receiver may keep arbitrarily
large quantum registers and act on them jointly; it cannot import them
coherently into a fresh acquisition.  Preparations and controls are one
common rule on every formal history, including null histories.  The
conditional old receiver states need not coincide.

For a history-dependent policy $\pi$, hard pathwise reservations define
$N_\pi=\sup_\zeta\sum_{h\in\zeta}n(h)$ and
$Q_\pi=\sup_\zeta\sum_{h\in\zeta}n(h)^2$.
Theorem~\ref{thm:quadraticbudget89} proves that the optimum over
$N_\pi\le N$, $Q_\pi\le Q$, $N\le Q\le N^2$, has coherent order
$\min\{1,\sqrt Qs\}$; its other two rows are unchanged.
The constants and the interval $0<s\le s_0$ depend on the fixed
$E,H,\Lambda$, but not on the profile, budget or unknown remainder.
A common bounded readout gives the lower without learning that remainder.
The individual-policy upper retains both $\alpha\sqrt{Q_\pi}$ and
$\sqrt{rN_\pi}$; a large cost is not a lower information certificate
for a policy which erases its output.

The second result concerns exact success, rather than local orders.
For every $d\ge2$, a finite family of ordered bases with simultaneous
Haar second moments exists.  Compare the scalar measurement $C_y=I/d$
with a once-chosen member $E_y^b(t)=(1-t)I/d+tP_y^b$, using that same
device twice.  The decision is scalar versus the family.  Put
$L=2(d+1)-t^2$ and assign these alternatives priors
$(d+1-t^2)/L$ and $(d+1)/L$.  Theorem~\ref{thm:operationalhierarchy90}
gives the exact values
\begin{equation}\label{eq:fronthierarchy90}
\begin{split}
 P_{\rm ca}&=(d+1)/L,\\
 P_{\rm reset}&=(d+1)/L+t^2(d-1)/(2dL),\\
 P_{\rm all}&=(d+1+t^2)/L.
\end{split}
\end{equation}
Here the first class completely measures and reprepares between calls;
the second allows arbitrary unit-reset receiver feedback; the third
allows all adaptive controls.  For $t>0$ the inequalities are strict,
including the full-rank interval $0<t<1$.  Two independent maximally
entangled pairs attain the middle value; an antisymmetric two-input
state attains the last.  A swap-filter spectral identity bounds every
reset policy, including all intervening receiver instruments.
For the scalar qubit measurement and the six ordered Pauli measurements,
the exact values are $3/5$, $13/20$, and $4/5$.
This finite ensemble result neither asserts a gap for every fixed pair
nor identifies a coherent-memory-dimension hierarchy.

\subsection{A general optimization principle and equal priors}
Theorem~\ref{thm:resetvariational91} applies to every finite two-call
ordered-measurement decision experiment, with arbitrary priors and
rewards.  If $g_y(\rho)$ is the best fresh-reference value after a first
label without receiver feedback and $c_y$ is its upper concave envelope
on all density matrices, then
\[
 P_{\rm reset}=\max_{\rho}\sum_yc_y(\rho).
\]
A common first marginal is essential.  At most
$(\operatorname{rank}\rho)^2$ instrument outcomes per first label
suffice, and the quantum receiver has dimension at most $d_1d_2$.
Theorem~\ref{thm:resetdual91} supplies exact Hermitian majorants and
contact conditions.  Mixed density-matrix atoms cannot be replaced
silently by pure states.  This is a finite-dimensional characterization
of optimal values, not a polynomial-time optimization claim.

For the same finite basis devices, Theorem~\ref{thm:equalprior91}
removes the special prior in \eqref{eq:fronthierarchy90}.  With equal
hypothesis priors its exact values are
\[
 \frac12+\frac{t^2(d-1)}{2d(d+1)},\qquad
 \frac12+\frac{t^2(d-1)}{2d^2},\qquad
 \frac12+\frac{t^2}{2d}.
\]
For the seven qubit devices at $t=1$ these become
$7/12<5/8<3/4$.  The middle upper permits all receiver feedback,
but its attaining strategy uses none.  The alternative index is chosen
once for both calls; independent redraws give no signal.  The explicit
qubit calculation and a channel/prior/law stability estimate follow the
theorem.

Theorem~\ref{thm:classicalroof91} identifies the classical optimum by
restricting the envelope atoms to pure states.  Thus two explicit
convexifications distinguish complete classical readout from quantum
receiver storage in every finite experiment.
Corollary~\ref{cor:gapcertificate91} gives necessary and sufficient
primal--dual certificates for a strict gap, and contact certificates
for equality.  Their majorization inequalities are universal; no
sampling procedure is asserted to certify them.

These statements share one operational interface and one precise
fresh-preparation cut, but have different quantifiers.  The local profile
law classifies shrinking fixed-tube signals at arbitrary hard resources.
The variational principle treats exact two-call values for arbitrary
finite experiments.  The two basis theorems solve particular experiments
inside that general class.  Neither exact calculation supplies uniform
constants for the local theorem, and the local theorem is not a premise
of either exact score proof.

\subsection{Geometry underlying the allocation law}
The positive covariance operator $\mathcal C_E$ acts on zero-sum
Hermitian tuples and incorporates the coupled measurement normalization.
For $M=(E+F)/2$ and $H=F-E$, the global certificate is
\begin{equation}\label{eq:frontupper84}
 D_N(E,F)\le\min\left\{2,2\sqrt{k+1}
 \sqrt{N\langle H,(\mathcal C_M+N^{-1}\operatorname{Id})^{-1}H\rangle}
 \right\}.
\end{equation}
Theorem~\ref{thm:supportkernel84} identifies its complete kernel in
support coordinates.  This is an upper certificate for arbitrary
pairs, not an asserted global matching metric.  On a nonstationary
fixed unitary orbit $E_j(t)=e^{itG}E_je^{-itG}$, the finite-angle theorem is
\begin{equation}\label{eq:frontorbit84}
 D_N(E,E(t))\asymp_{E,G}
 \begin{cases}
 \min\{1,\sqrt N|t|\},&G\in\mathcal W_E,\\
 \min\{1,N|t|\},&G\notin\mathcal W_E.
 \end{cases}
\end{equation}
The support span is the Hermitian Kraus-product span.  The channel
metrological criterion and correction principle are credited to
Zhou--Jiang \cite{ZJ84}; the finite-pair estimates retain explicit
quadratic channel errors.  Normalized tangent realizations extend
these estimates to the whole legal tube, including unknown remainders
and one-sided support openings.

The width-cap consequence, for $1\le b\le N$, is
\begin{equation}\label{eq:frontwidth87}
 D_N^{[b]}(E,F)\asymp_{E,H,\Lambda}
 \begin{cases}
 \min\{1,\sqrt Ns\},&H\in\operatorname{Ran}\mathcal C_E,\\
 \min\{1,\sqrt{Nb}\,s\},&J_j=0\ (\forall j),
                          \Gamma_E(H)\notin\mathcal W_E,\\
 \min\{1,Ns\},&J_j\ne0\ (\exists j).
 \end{cases}
\end{equation}
The complete profile refines this cap through $V$, rather than through
its largest group alone.  The upper uses conditional fidelity, not
unconditional tensorization after feedback.  The lower first pays the
full $m_r\kappa s^2$ block errors and then combines unequal signals.
These distinctions are essential to uniformity in the unknown remainder.

The variational and equality statements also have quantitative forms.
Corollary~\ref{cor:contactdefect91} decomposes a certified reward loss
into nonnegative spectral, majorization, and leaf-decision defects.
Lemma~\ref{lem:stableswap91} and Corollary~\ref{cor:resetstability91}
control, at the relative scale $\varepsilon/t^2$, the deviation of
near-optimal reset Gram matrices from scalar matrices.  This concerns
the selected normal form, not unique dilations or physical calibration.

\subsection{Notation, organization and scope}
For positive trace-class $R,S$, $f(R,S)=\|\sqrt R\sqrt S\|_1$ is root
fidelity.  For states, $d_B(R,S)^2=2(1-f(R,S))$; the inherited symbol
$\beta$ for a state Bures distance denotes this same distance when used
with two state arguments.  Named coefficients such as $\beta_h$ are
bound coefficients, not an additional distance.  $D(\pi;E,F)$ denotes
one policy's separation; a resource subscript denotes a supremum.
$V/N$ is an effective allocation width, not an integer or memory dimension.

The proof proceeds from covariance and its support kernel to normalized
curves, tangent neighborhoods, complete groups, conditional reset
accumulation, and hard quadratic budgets.  Section~\ref{sec:operationalmemory90}
then gives the exact finite hierarchy and its stability.  Balanced-interior
geometry, coding, learning and affine repair are retained with their full
proofs in the current linked supplement.  The complete research edition
also retains every earlier mathematical section.  The independent
structural article is not a premise of either main result.

The fixed-tube constants may deteriorate at changing support floors,
vanishing tangents or closing correction gaps.  Neither local theorem
is a uniform full-boundary metric, entropy or learning theorem.  The
finite hierarchy supplies a different, exact, all-dimension statement on
its specified ensemble family.  The historical analytic programme and
independent specialist priority assessment are separate obligations.
Our comparison with the memory-dimension hierarchy of Zonnios--Binder
\cite{ZB90} and with accessible post-measurement states \cite{EQ90}
keeps these interfaces and resources distinct.
